% ID: SEL-D0064-OB-E04
% ESTADO: cotejada. Segunda pasada visual y borrador compilado revisados.
% Fuente: Selectividad/sel_2010_matematicas_aplicadas/mat_apli/mat_apli_sept/mat_apli_sept/EXAMEN-ANDALUCÍA-3.pdf
% SHA-256: 9c3d2d628115a2f20f17b47171106ddd90f214b5f6804bdf9a7dcb7037c75efd
% Página: 2; opción B; ejercicio original 4.

\begin{enumerate}[label=\alph*)]
\item \textbf{(1.25 puntos)} La altura de los alumnos de una Universidad sigue una distribución Normal de media desconocida y desviación típica 11 cm. Calcule el tamaño mínimo que ha de tener una muestra aleatoria de esos alumnos para que el error cometido al estimar la altura media sea inferior a 1 cm, con un nivel de confianza del $98\%$.
\item \textbf{(1.25 puntos)} Dada la población $\{10,12,17\}$, escriba todas las muestras de tamaño 2 mediante muestreo aleatorio simple y calcule la media y la desviación típica de las medias muestrales.
\end{enumerate}
